% GENERATED FILE. Edit canonical lessons, then run npm run generate:books.
% canonical-source-hash: fnv1a64:9d8185ed9113b85b
% canonical-lessons: GE-C329-anrichten, GE-C329-garnieren, GE-C329-servieren, GE-C329-abkuehlen, GE-C329-knusprig, GE-R329-review-326-329, GE-R329-review-final-329

\chapter{Guten Appetit}
\label{ch:ge-guten-appetit}
\hlchaptermodality{329}

\begin{chapteropening}
\textbf{By the end of this chapter:} \emph{I can describe serving a meal: to arrange on the plate, to garnish, to serve, to cool and crispy.}

Complete the last lesson of chapter 329: I can describe serving a meal: to arrange on the plate, to garnish, to serve, to cool and crispy.
\end{chapteropening}

\section[anrichten]{anrichten — to arrange on the plate; to cause (damage)}

\label{lesson:GE-C329-anrichten}



\begin{quote}
Before the new one: say the German for to taste for seasoning, then the German for to season.
\end{quote}

\subsection*{You'll want to know: anrichten}
\textbf{anrichten} — \textquotedblleft{}to arrange on the plate; to cause (damage)\textquotedblright{}.

\begin{quote}
Sie richtete das Essen auf großen Tellern an.
\end{quote}

\begin{quote}
\emph{She arranged the food on large plates.}
\end{quote}

Separable. Also to cause damage: \emph{Schaden anrichten}.

\begin{cousinweb}
\emph{anrichten} is built on \emph{richten} \textquotedblleft{}to put right, to prepare\textquotedblright{}, Old High German \emph{rihten}, Proto-Germanic \emph{rihtijaną} (EtymDB).
\end{cousinweb}

\begin{grammarlens}[title={What you've built}]
One more word for telling, explaining and describing at B1.
\end{grammarlens}

\subsection*{Your turn}
\begin{itemize}
\raggedright
  \item \emph{Say it:} \emph{anrichten}
  \item \emph{Say it:} \emph{anrichten}, once more
  \item \emph{Say it:} a sentence of your own with \emph{anrichten}
  \item \emph{Recall:} say the German for to taste for seasoning, then the German for to season, then say \emph{anrichten} again
\end{itemize}

\begin{tcolorbox}[breakable,colback=teal!4,colframe=teal!35!black,title={Before you move on}]
What does anrichten mean? (\textquotedblleft{}To arrange on the plate; to cause (damage)\textquotedblright{}.) Say it once more.
\end{tcolorbox}
\section[garnieren]{garnieren — to garnish}

\label{lesson:GE-C329-garnieren}



\begin{quote}
Before the new one: say the German for to season, then the German for to arrange on the plate.
\end{quote}

\subsection*{You'll want to know: garnieren}
\textbf{garnieren} — \textquotedblleft{}to garnish\textquotedblright{}.

\begin{quote}
Er garnierte den Fisch mit Zitrone.
\end{quote}

\begin{quote}
\emph{He garnished the fish with lemon.}
\end{quote}

Regular, without \emph{ge-}: \emph{hat garniert}.

\begin{cousinweb}
\emph{garnieren} was borrowed from French \emph{garnir} \textquotedblleft{}to furnish, to equip\textquotedblright{}, which is usually traced to a Frankish word; EtymDB lists no parent.
\end{cousinweb}

\begin{grammarlens}[title={What you've built}]
One more word for telling, explaining and describing at B1.
\end{grammarlens}

\subsection*{Your turn}
\begin{itemize}
\raggedright
  \item \emph{Say it:} \emph{garnieren}
  \item \emph{Say it:} \emph{garnieren}, once more
  \item \emph{Say it:} a sentence of your own with \emph{garnieren}
  \item \emph{Recall:} say the German for to season, then the German for to arrange on the plate, then say \emph{garnieren} again
\end{itemize}

\begin{tcolorbox}[breakable,colback=teal!4,colframe=teal!35!black,title={Before you move on}]
What does garnieren mean? (\textquotedblleft{}To garnish\textquotedblright{}.) Say it once more.
\end{tcolorbox}
\section[servieren]{servieren — to serve (food)}

\label{lesson:GE-C329-servieren}



\begin{quote}
Before the new one: say the German for to arrange on the plate, then the German for to garnish.
\end{quote}

\subsection*{You'll want to know: servieren}
\textbf{servieren} — \textquotedblleft{}to serve (food)\textquotedblright{}.

\begin{quote}
Der Kellner servierte die Suppe heiß.
\end{quote}

\begin{quote}
\emph{The waiter served the soup hot.}
\end{quote}

Regular, without \emph{ge-}: \emph{hat serviert}.

\begin{cousinweb}
\emph{servieren} was borrowed from French \emph{servir}, from Latin \emph{servīre} \textquotedblleft{}to serve\textquotedblright{} (EtymDB).
\end{cousinweb}

\begin{grammarlens}[title={What you've built}]
One more word for telling, explaining and describing at B1.
\end{grammarlens}

\subsection*{Your turn}
\begin{itemize}
\raggedright
  \item \emph{Say it:} \emph{servieren}
  \item \emph{Say it:} \emph{servieren}, once more
  \item \emph{Say it:} a sentence of your own with \emph{servieren}
  \item \emph{Recall:} say the German for to arrange on the plate, then the German for to garnish, then say \emph{servieren} again
\end{itemize}

\begin{tcolorbox}[breakable,colback=teal!4,colframe=teal!35!black,title={Before you move on}]
What does servieren mean? (\textquotedblleft{}To serve (food)\textquotedblright{}.) Say it once more.
\end{tcolorbox}
\section[abkühlen]{abkühlen — to cool down}

\label{lesson:GE-C329-abkuehlen}



\begin{quote}
Before the new one: say the German for to garnish, then the German for to serve.
\end{quote}

\subsection*{You'll want to know: abkühlen}
\textbf{abkühlen} — \textquotedblleft{}to cool down\textquotedblright{}.

\begin{quote}
Der Kuchen kühlte am Fenster ab.
\end{quote}

\begin{quote}
\emph{The cake cooled down by the window.}
\end{quote}

Separable. Perfekt with \emph{sein} when the thing itself cools: \emph{Der Kuchen ist abgekühlt}.

\begin{cousinweb}
\emph{abkühlen} joins \emph{ab} and \emph{kühlen} \textquotedblleft{}to cool\textquotedblright{} (EtymDB), an inherited Germanic verb and a cousin of English \emph{cool}.
\end{cousinweb}

\begin{grammarlens}[title={What you've built}]
One more word for telling, explaining and describing at B1.
\end{grammarlens}

\subsection*{Your turn}
\begin{itemize}
\raggedright
  \item \emph{Say it:} \emph{abkühlen}
  \item \emph{Say it:} \emph{abkühlen}, once more
  \item \emph{Say it:} a sentence of your own with \emph{abkühlen}
  \item \emph{Recall:} say the German for to garnish, then the German for to serve, then say \emph{abkühlen} again
\end{itemize}

\begin{tcolorbox}[breakable,colback=teal!4,colframe=teal!35!black,title={Before you move on}]
What does abkühlen mean? (\textquotedblleft{}To cool down\textquotedblright{}.) Say it once more.
\end{tcolorbox}
\section[knusprig]{knusprig — crispy, crunchy}

\label{lesson:GE-C329-knusprig}



\begin{quote}
Before the new one: say the German for to serve, then the German for to cool down.
\end{quote}

\subsection*{You'll want to know: knusprig}
\textbf{knusprig} — \textquotedblleft{}crispy, crunchy\textquotedblright{}.

\begin{quote}
Das Brot war außen knusprig und innen weich.
\end{quote}

\begin{quote}
\emph{The bread was crispy outside and soft inside.}
\end{quote}

Built on \emph{knuspern}, to nibble something crunchy.

\begin{cousinweb}
\emph{knusprig} is built on \emph{knuspern} \textquotedblleft{}to nibble crunchily\textquotedblright{}, a word that imitates the sound; EtymDB lists no parent.
\end{cousinweb}

\begin{grammarlens}[title={What you've built}]
One more word for telling, explaining and describing at B1.
\end{grammarlens}

\subsection*{Your turn}
\begin{itemize}
\raggedright
  \item \emph{Say it:} \emph{knusprig}
  \item \emph{Say it:} \emph{knusprig}, once more
  \item \emph{Say it:} a sentence of your own with \emph{knusprig}
  \item \emph{Recall:} say the German for to serve, then the German for to cool down, then say \emph{knusprig} again
\end{itemize}

\begin{tcolorbox}[breakable,colback=teal!4,colframe=teal!35!black,title={Before you move on}]
What does knusprig mean? (\textquotedblleft{}Crispy, crunchy\textquotedblright{}.) Say it once more.
\end{tcolorbox}
\section[(dialogue)]{Noch einmal alles}

\label{lesson:GE-R329-review-326-329}



\begin{quote}
Say the German for to preheat and to grease.
\end{quote}

\subsection*{Your turn}
Say these aloud:

\begin{itemize}
\raggedright
  \item to preheat, to grease, the baking tray, the baking dish, to brown
  \item to simmer, to bring to the boil, to steam, to brown quickly, the lid
  \item the sieve, to drain, to drip dry, to taste for seasoning, to season
  \item to arrange on the plate, to garnish, to serve, to cool down, crispy
\end{itemize}

\begin{tcolorbox}[breakable,colback=teal!4,colframe=teal!35!black,title={Before you move on}]
To preheat? (\textbf{vorheizen}.) To grease? (\textbf{einfetten}.) The baking tray? (\textbf{das Backblech}.) The baking dish? (\textbf{die Auflaufform}.) To brown? (\textbf{bräunen}.) To simmer? (\textbf{köcheln}.) To bring to the boil? (\textbf{aufkochen}.) To steam? (\textbf{dünsten}.) To brown quickly? (\textbf{anbraten}.) The lid? (\textbf{der Deckel}.) The sieve? (\textbf{das Sieb}.) To drain? (\textbf{abgießen}.) To drip dry? (\textbf{abtropfen}.) To taste for seasoning? (\textbf{abschmecken}.) To season? (\textbf{würzen}.) To arrange on the plate? (\textbf{anrichten}.) To garnish? (\textbf{garnieren}.) To serve? (\textbf{servieren}.) To cool down? (\textbf{abkühlen}.) Crispy? (\textbf{knusprig}.)
\end{tcolorbox}
\section[(dialogue)]{Guten Appetit — once more}

\label{lesson:GE-R329-review-final-329}



\begin{quote}
Say the German for to arrange on the plate and to garnish.
\end{quote}

\subsection*{Your turn}
\begin{itemize}
\raggedright
  \item \emph{Say it:} to arrange on the plate, to garnish, to serve, to cool down, crispy
\end{itemize}

\begin{tcolorbox}[breakable,colback=teal!4,colframe=teal!35!black,title={Before you move on}]
To arrange on the plate? (\textbf{anrichten}.) To garnish? (\textbf{garnieren}.) To serve? (\textbf{servieren}.) To cool down? (\textbf{abkühlen}.) Crispy? (\textbf{knusprig}.)
\end{tcolorbox}
